\section{За пределами решёток: мозаичные раскраски и пол метода}
\label{sec:tilings}

Все конструкции этой работы, от раскраски решёткой общего положения до
ламинирования и продуктового исчисления, устроены одинаково: пространство
режется на ячейки Вороного некоторой решётки $\Lambda$, и ячейка
$V_0+u$ получает цвет $[u]\in\Lambda/\Gamma$. Естественно спросить, чего стоит
эта решёточность --- и сколько вообще можно выиграть, если от неё отказаться.
В плоскости у этого вопроса есть собственная литература: нижние границы
интервальной задачи получают из $\varepsilon$-графов единичных расстояний
\cite{Exoo2005,ChJW,Parts2023}, а верхние --- мозаиками, от неправильных
шестиугольников \cite{deGreyParts} до плиток разных форм \cite{Parts2023}.
Здесь мы смотрим на тот же вопрос в произвольной размерности.

Мы (i) строим строго более широкий класс раскрасок, в котором решётка
$\Lambda$ отсутствует вовсе, (ii) доказываем, что весь этот класс упирается в
пол $2^n$ цветов, (iii) показываем структурную причину, по которой
дополнительная свобода не окупается в \emph{чемпионских} индексах, и
(iv) показываем, где она окупается: на индексах, запрещённых для решёточной
схемы арифметикой, мозаики строго обгоняют её --- и мы предъявляем для этого
точные сертификаты (\S\ref{ssec:quant}). Итог точнее расхожего <<решёточность
не мешает>>: она не мешает ровно там, где живут рекорды --- на индексах-нормах
--- и строго мешает между ними; ограничение сверху --- сам мозаичный принцип.

\subsection{Степенные (лагерровы) мозаики}
\label{ssec:power}

\begin{definition}[раскраска мозаикой]\label{def:tiling}
Пусть $\Gamma\subset\R^n$ --- решётка, $t_1,\dots,t_k\in\R^n$ --- узлы и
$w_1,\dots,w_k\in\R$ --- веса. Придадим узлу $t_i+g$ ($g\in\Gamma$) вес $w_i$ и
возьмём \emph{степенную} (лагеррову) диаграмму этого $\Gamma$-периодического
набора: ячейка узла $s$ с весом $w_s$ есть
$\{x:\ |x-s|^2-w_s\le|x-s'|^2-w_{s'}\ \forall s'\}$. Ячейки выпуклы, замощают
$\R^n$, ячейка узла $t_i+g$ есть $P_i+g$; красим ячейку в цвет $i$.
\end{definition}

Раскраска использует $k$ цветов (узел с пустой ячейкой цвета не тратит). Две
точки одного цвета лежат либо в одной ячейке, либо в $P_i+g$ и $P_i+g'$; обе
ячейки выпуклы, поэтому реализуемые ими расстояния заполняют отрезок
$[\dist(g-g',P_i-P_i),\ \cdot\,]$. Значит, положив
\[
\Delta=\max_i\diam P_i,\qquad
\mathrm{gap}=\min_i\ \min_{0\ne g\in\Gamma}\dist(g,\ P_i-P_i),
\]
получаем: раскраска не реализует ни одного расстояния из $(\Delta,\mathrm{gap})$,
то есть после нормировки
\begin{equation}\label{eq:tilewidth}
\chi\bigl(\R^n,[1,d]\bigr)\le k,\qquad d:=\mathrm{gap}/\Delta\ \ge 1 .
\end{equation}

\begin{proposition}[схема Вороного --- частный случай]\label{prop:contains}
Пусть $\Gamma\subseteq\Lambda$ индекса $k$, узлы $t_i$ --- полная система
представителей $\Lambda/\Gamma$, все $w_i=0$. Тогда степенная диаграмма есть
диаграмма Вороного решётки $\Lambda$, все $P_i$ --- сдвиги $V_0$, и
\[
P_i-P_i=2V_0,\qquad \dist(g,2V_0)=2\dist(g/2,V_0)=D(g),
\]
так что \eqref{eq:tilewidth} превращается ровно в $d=D_{\min}/\diam V_0$.
\end{proposition}

Таким образом определение~\ref{def:tiling} содержит всю работу как срез
$w=0$, $T=\Lambda/\Gamma$, а новая свобода --- непрерывная и заведомо
нерешёточная: узлы не обязаны образовывать группу, а ячейки одного цвета не
обязаны быть конгруэнтны ячейкам другого.

Реализация --- \texttt{core/power\_coloring.py}. Величина $\mathrm{gap}$
считается алгоритмом Вульфа для задачи о точке минимальной нормы в политопе: линейный
оракул над (никогда не выписываемым) множеством $\{g-a+b\}$ разностей вершин
даётся парой $\arg\max/\arg\min$ скалярного произведения, а каждая итерация
предъявляет разделяющий слой, дающий \emph{нижнюю} оценку
$\dist(g,P-P)\ge\langle g,e\rangle-\operatorname{width}(P,e)$. Поэтому
сообщаемое $d$ никогда не завышено. Контроль: для
$A_2/7$, $A_2/9$, $\Z^3/27$, $A_3^*/27$, $D_4/81$, $A_4^*/81$ воспроизводятся
замкнутые формы $\sqrt7/2$, $\sqrt3$, $2/\sqrt3$, $2/\sqrt{5/3}$, $\sqrt2$,
$\sqrt2$ с ошибкой $<10^{-9}$, а для чемпионов работы --- $d=1{,}0163393$ при
$k=45$ и $d=0{,}9903509$ при $k=44$ (совпадение до последнего знака с
\S\ref{sec:comp} и \S\ref{sec:open}).

\subsection{Пол мозаичного принципа}
\label{ssec:floor}

\begin{theorem}[пол $2^n$]\label{thm:floor}
Пусть $\R^n$ замощён выпуклыми плитками диаметра $<1$, замощение
$\Gamma$-периодично, плитки одного цвета попарно удалены больше чем на~$1$, а
$P_1,\dots,P_k$ --- плитки одного периода. Тогда
\[
k\ \ge\ \bigl(1+\theta^{1/n}\bigr)^n\ \ge\ 2^n,
\qquad
\theta:=\frac{\omega_n2^{-n}}{\max_i\vol P_i}\ \ge\ 1 .
\]
Равенство потребовало бы, чтобы плитка была шаром, поэтому оно недостижимо.
\end{theorem}

\begin{proof}
Зафиксируем цвет $i$. Тела $P_i+g+B(0,\tfrac12)$, $g\in\Gamma$, попарно не
пересекаются (иначе нашлись бы две точки цвета $i$ на расстоянии ровно~1),
поэтому $\det\Gamma\ge\vol\bigl(P_i\oplus B(0,\tfrac12)\bigr)$ для каждого $i$.
Плитки точно заполняют период: $\det\Gamma=\sum_i\vol P_i$. Пусть
$V=\max_i\vol P_i$ достигается на $j$. Неравенство Брунна--Минковского даёт
\[
\det\Gamma\ \ge\ \Bigl(V^{1/n}+\bigl(\omega_n2^{-n}\bigr)^{1/n}\Bigr)^{n}
= \bigl(1+\theta^{1/n}\bigr)^{n}\,V,
\]
а изодиаметрическое неравенство при $\diam P_j\le1$ даёт
$V\le\omega_n2^{-n}$, то есть $\theta\ge1$. Вместе с
$\det\Gamma=\sum_i\vol P_i\le kV$ это и есть утверждение.
\end{proof}

Теорема~\ref{thm:floor} не знает ни о какой решётке $\Lambda$ и потому
применима к определению~\ref{def:tiling} целиком. Для решёточных раскрасок
известен более сильный пол Кулсона $2^{n+1}-1$, а в плоскости для мозаичных
раскрасок --- оценка $\ge6$ \cite{Coulson2004}; сила теоремы~\ref{thm:floor}
--- в произвольной размерности. В частности, \emph{никакое} обобщение
мозаичного типа не даст асимптотики лучше $2^n$, тогда как классическая
оценка Лармана--Роджерса есть $3^n$: весь запас метода --- это множитель
$(3/2)^n$.

\subsection{Почему нерешёточная свобода не окупается}
\label{ssec:whynot}

\begin{proposition}[редукция к одной форме]\label{prop:shape}
Зафиксируем решётку периодов $\Gamma$ и нормируем запрещённое расстояние
единицей. Назовём выпуклое тело $P$ \emph{допустимым}, если $\diam P\le1$ и
$\dist(g,P-P)\ge1$ для всех $g\in\Gamma\setminus\{0\}$. Это условие не
упоминает ни остальные плитки, ни цвет, и наследуется подмножествами. Поэтому
для любой раскраски по определению~\ref{def:tiling} с решёткой периодов
$\Gamma$
\[
k\ =\ \frac{\det\Gamma}{\text{средний объём плитки}}\ \ge\
\frac{\det\Gamma}{\vol P^*},\qquad
P^*:=\arg\max\{\vol P:\ P\ \text{допустимо}\},
\]
и всякая раскраска, достигающая границы, состоит из плиток максимального
объёма, то есть из копий одной формы.
\end{proposition}

Отсюда --- структурный вывод. При фиксированной $\Gamma$ все $k$ плиток
подчинены \emph{одному и тому же} условию допустимости, а их суммарный объём
закреплён. Значит, оптимум хочет сделать все плитки одинаковыми --- и это ровно
ситуация Вороного. Дополнительная свобода степенной диаграммы может окупиться
только через \emph{реализуемость замощения} (когда тело $P^*$ само по себе
решёткой не замощает), но никогда --- через условие на отдельную плитку.

Численно это подтверждается точно. Оптимизатор (CMA-ES по всем параметрам:
решётка периодов, узлы, веса --- от 21 до 230 переменных) сажался \emph{ровно} в
известный оптимум и получал всю дополнительную свободу:

\begin{center}
\begin{tabular}{@{}lrrl@{}}
\toprule
случай & $k$ & параметров & итог\\
\midrule
$\R^2$, шестиугольная $A_2$ & $7$  & $21$  & $1{,}3228757\to$ без улучшения\\
$\R^3$, ОЦК (оптимум Кулсона) & $15$ & $62$  & $1{,}0000000\to$ без улучшения\\
$\R^4$, чемпион \S\ref{sec:comp} & $45$ & $230$ & $1{,}0163393\to$ без улучшения\\
$\R^4$, лучшее при $k=44$ & $44$ & $225$ & $0{,}9903509\to$ без улучшения\\
\bottomrule
\end{tabular}
\end{center}

Калибровка силы оптимизатора: тот же CMA-ES со случайных стартов доходит в
$\R^2$ при $k=7$ лишь до $1{,}2396$ (истина $\sqrt7/2=1{,}3229$), то есть
отрицательный результат со случайных стартов был бы слабым свидетельством;
здесь же старт --- сам чемпион, и утверждение локальное и потому надёжное.
Механизм локальности виден из предложения~\ref{prop:shape}: в точке Вороного
все $k$ плиток конгруэнтны, поэтому \emph{все} активные ограничения (максимум
диаметра и минимум зазора) достигаются одновременно на каждой плитке; любое
шевеление весов увеличивает объём части плиток, а с ним --- их диаметр, и
уменьшает их зазор.

\subsection{Арифметическая квантизация: где свобода окупается}
\label{ssec:quant}

Предложение~\ref{prop:shape} объясняет, почему свобода не окупается в
\emph{чемпионских} индексах, но оставляет лазейку --- реализуемость
замощения, --- и она срабатывает. Решёточная схема реализует в плоскости
индекс $k$ подобной (то есть правильной шестиугольной) подрешёткой, а такая
существует, лишь когда $k$ --- норма кольца $\Z[\omega]$ (числа Лёшиана
$1,3,4,7,9,12,13,16,19,21,\dots$); на всех остальных $k$ решёточная лестница
ширин проваливается. Препятствие \emph{арифметическое}, а не геометрическое,
и мозаичная свобода определения~\ref{def:tiling} снимает именно его. В
плоскости этот эффект известен: де~Грей и Партс \cite{deGreyParts} обгоняют
решёточную схему неправильными шестиугольниками ровно на нелёшиановых $k$, а
Партс \cite{Parts2023} плитками разных форм доводит рекорды до
$d\approx1{,}4442$ при $k=8$ и $d\approx2{,}3470$ при $k=15$ (численно, без
сертификатов).

Мы воспроизводим этот эффект внутри класса~\ref{def:tiling} в
\emph{сертифицированной} форме: конфигурация из $N$ орбит
$P=\bigcup_i(\Gamma+t_i)$ с лагерровыми весами, найденная CMA-ES, округляется
до рациональной, после чего вершины ячеек, замощение и все разделения
проверяются в точной арифметике дробей --- без единого извлечения корня
(\texttt{core/semilattice\_cert.py}; валидация --- семь классических ступеней,
воспроизведённых точными значениями $d^2$ от $3/4$ до $37/4$).

\begin{keybox}
\textbf{[С]}\quad
$\chi\bigl(\R^2,[1,\ell]\bigr)\le8$ при $\ell\le1{,}441350$,\qquad
$\chi\bigl(\R^2,[1,\ell]\bigr)\le15$ при $\ell\le2{,}281534$
\end{keybox}

\noindent против $1{,}400000$ ($=7/5$) и $2{,}186607$ лучшей решёточной схемы
при тех же $N$ (исчерпывающий перебор форм Эрмита и форм плоских решёток).
Компактные сертификаты со знаменателем $20$ дают $d=1{,}424440$ и
$d=2{,}219896$ --- слабее, но проверяемые вручную. Ячейки заметно
нерешёточные: при $N=8$ это пяти-, шести- и семиугольники с выровненными
диаметрами. Насколько нам известно, это первые \emph{точно
сертифицированные} нерешёточные мозаичные раскраски плоскости; численные
рекорды Партса чуть сильнее, и довести их конфигурации до точных
сертификатов --- прямая задача для того же сертификатора. Проверка тройная:
независимое перечисление вершин, оценщик \S\ref{ssec:power} (совпадение до
$10^{-9}$) и Монте-Карло прямо по определению раскраски
(\texttt{results/semilattice\_cert\_N\{8,15\}.json}).

Отрицательная сторона той же монеты важнее для больших размерностей: узкое
место оценки $\chi(\R^{10})\le45619$ --- плоский проставок, которому нужно
$d\ge\sqrt7$ при минимальном числе цветов. Геометрия разрешает $k=17$
(инрадиусный экран, $16{,}44$), арифметика решёточной схемы требует $k=19$
(\S\ref{ssec:planefloor}). Мозаичная свобода могла бы отыграть $17$ или $18$
--- но не отыгрывает: минимум по $18$-точечным конфигурациям максимума
диаметра ячейки равен $0{,}285366$ при пределе
$1/(\sqrt3/2+\sqrt7)=0{,}284756$ --- промах $0{,}21\,\%$ и потолок
$d\lesssim2{,}638$ при нужных $2{,}6458$; поиск ширины устойчиво упирается в
$2{,}598076$ --- значение решётки $k=16$ с лишними точками. Здесь квантизация,
по-видимому, не связывает: $19$ --- минимум и для мозаик (статус
\textbf{[Ч]}: промах измерен, но не доказан), а оценки $45619$/$28812$
устойчивы к нерешёточной свободе.

\subsection{Точный учёт метода и ландшафт запаса}
\label{ssec:accounting}

Условие допустимости решёточной раскраски $D(v)\ge2R$ для $v\in\Gamma$ есть в
точности $\Gamma\cap\operatorname{int}2C=\{0\}$ для \emph{округлённой ячейки}
$C=V_0\oplus RB$, то есть $\Gamma$ --- упаковочная решётка тела $C$. Отсюда
тождество
\begin{equation}\label{eq:accounting}
k=\frac{F(\Lambda)}{\delta_C(\Gamma)},\qquad
F(\Lambda)=\frac{\vol(V_0\oplus RB)}{\det\Lambda}\ \ge\ \bigl(1+\Theta^{1/n}\bigr)^n,
\end{equation}
где $\Theta$ --- плотность покрытия $\Lambda$, а $\delta_C\le1$ --- плотность
упаковки телом $C$. Читается это так: \emph{родитель должен быть хорошим
покрытием, подрешётка --- хорошей упаковкой округлённой ячейки}. Все
конструкции работы берут $\Gamma=\alpha\Lambda$, подобную $\Lambda$, и тем
самым требуют от одной решётки обоих качеств сразу.

Таблица~\ref{tab:landscape} собирает запас каждого рекорда против обоих полов
(\texttt{core/paradigm\_floor.py}, артефакт
\texttt{results/power\_tilings.json}). Столбец $k^{1/n}$ --- эффективное
основание.

\begin{table}[htbp]
\centering
\caption{Ландшафт запаса. $\Theta$ выведена из точных
$(\lambda_1^2,R^2,\det)$ и сверена с независимыми замкнутыми формами
$2\pi/(3\sqrt3)$, $\pi^4/24$, $\pi^{12}2^{12}/12!$ до $10^{-12}$.}
\label{tab:landscape}
\begin{tabular}{@{}rrrlrrrr@{}}
\toprule
$n$ & цветов & $k^{1/n}$ & родитель & $\rho$ & $\Theta$ & $k/F_{\min}$ & $k/2^n$\\
\midrule
$2$  & $7$          & $2{,}646$ & $A_2$          & $1{,}155$ & $1{,}21$ & $1{,}6$ & $1{,}8$\\
$3$  & $15$         & $2{,}466$ & $A_3^*$        & $1{,}291$ & $1{,}46$ & $1{,}5$ & $1{,}9$\\
$4$  & $45$         & $2{,}590$ & общего пол.    & ---       & ---      & ---     & $2{,}8$\\
$5$  & $132$        & $2{,}655$ & общего пол.    & ---       & ---      & ---     & $4{,}1$\\
$6$  & $343$        & $2{,}646$ & $E_6^*$        & $1{,}414$ & $2{,}65$ & $3{,}2$ & $5{,}4$\\
$7$  & $1029$       & $2{,}694$ & лам.\ $E_6^*$  & ---       & ---      & ---     & $8{,}0$\\
$8$  & $2401$       & $2{,}646$ & $E_8$          & $1{,}414$ & $4{,}06$ & $4{,}5$ & $9{,}4$\\
$9$  & $7203$       & $2{,}683$ & лам.\ $E_8$    & ---       & ---      & ---     & $14{,}1$\\
$10$ & $28812$      & $2{,}792$ & лам.\ $E_8{+}A_2$ & ---    & ---      & ---     & $28{,}1$\\
$11$ & $3^{11}$     & $3{,}000$ & башня          & $1{,}791$ & ---      & ---     & $86{,}5$\\
$12$ & $3^{12}$     & $3{,}000$ & $K_{12}$       & $1{,}633$ & $17{,}8$ & $28{,}2$& $129{,}7$\\
$24$ & $7^{12}$     & $2{,}646$ & $\Lambda_{24}$ & $1{,}414$ & $7{,}90$ & $287$   & $825$\\
\bottomrule
\end{tabular}
\end{table}

Столбец $k/F_{\min}$ --- это $1/\delta_C$ по \eqref{eq:accounting}, то есть
потеря на упаковке округлённой ячейки. Она нигде не мала и почти всюду близка к
$1/\delta_{\max}(n)$: при $n=3$ это $1{,}5$ против $1/0{,}7405=1{,}35$, при
$n=8$ --- $4{,}5$ против $3{,}94$, при $n=12$ --- $28{,}2$ против $20{,}2$. Иными
словами, оставшийся запас метода --- множитель порядка $1{,}4$--$2$ в
размерностях $10$--$12$, а не порядки; в размерностях $3$ и $8$ рекорды стоят в
$10$--$15\,\%$ от реалистичного пола.

\subsection{Куда ещё можно целиться}
\label{ssec:window}

Пол \eqref{eq:accounting} объёмный, но есть второй, обычно более сильный.
Инрадиусная лемма $D(v)\le|v|-\lambda_1(\Lambda)$ превращает допустимость
$D(v)>2R$ в $\lambda_1(\Gamma)>2R+\lambda_1(\Lambda)$; решётка с таким минимумом
есть упаковка шарами радиуса $\lambda_1(\Gamma)/2$, поэтому
\begin{equation}\label{eq:packfloor}
k\ \ge\ \frac{\omega_n\bigl(R+\lambda_1(\Lambda)/2\bigr)^{n}}
             {\delta_{\max}(n)\,\det\Lambda}.
\end{equation}

\begin{center}
\begin{tabular}{@{}lrrrrrr@{}}
\toprule
родитель & $n$ & $\lambda_1(\Gamma)\ge$ & \eqref{eq:packfloor} &
$(1+\Theta^{1/n})^n$ & рекорд & запас\\
\midrule
$A_3^*$ (ОЦК) & $3$  & $2{,}291$ & $11$      & $10$      & $15$      & $1{,}36$\\
$E_8$         & $8$  & $3{,}414$ & $1154$    & $532$     & $2401$    & $2{,}08$\\
$K_{12}$      & $12$ & $5{,}266$ & $111\,019$& $18\,826$ & $3^{12}$  & $4{,}79$\\
\bottomrule
\end{tabular}
\end{center}

Оценка \eqref{eq:packfloor} строга при $n=3$ и $n=8$ (плотнейшие упаковки
доказаны Хейлзом и Вязовской); при $n=12$ оптимальность $K_{12}$ открыта, так
что строка $K_{12}$ условна. Читается таблица так: подрешётка $\Gamma\subset
K_{12}$, улучшающая $3^{12}$, обязана иметь $\lambda_1^2\ge28$ и индекс в окне
$[111\,019,\ 3^{12}]$ --- запас почти в пять раз, самый большой из всех
родителей. Эйзенштейнова $(3+\omega)K_{12}$ с индексом $7^6=117\,649$ лежит
чуть выше пола, но выбывает по другой причине: у неё $D_{\min}$ равно ровно
планарному полу $\sqrt{7/3}\lambda_1=3{,}055$ при нужных $3{,}266$
(\S\ref{sec:product}). Лестница подобий в этом окне пуста (эйзенштейновы нормы
идут $7,9,12,13,\dots$), поэтому искать надо $\Z[\omega]$-подмодули $K_{12}$ с
произвольной нормой определителя --- это и есть направление, на которое
указывает учёт.

\subsection{Где кончается эйзенштейнова лестница}
\label{ssec:kl}

Для подобной подрешётки $\Gamma=\alpha\Lambda$ вся допустимость сводится к
одному неравенству на $\rho=\diam V_0/\lambda_1=2R/\lambda_1$: $\rho\le2$ при
$\alpha=3$ (это $3^n$ цветов) и $\rho\le\sqrt{7/3}=1{,}52753$ при
$\alpha=3+\omega$ (это $7^{n/2}=2{,}6458^n$). С другой стороны, прямо из
определений
\begin{equation}\label{eq:rhoid}
\frac{\Theta(\Lambda)}{\delta(\Lambda)}
=\Bigl(\frac{2R}{\lambda_1}\Bigr)^{n}=\rho^{n},
\end{equation}
а $\Theta\ge1$, поэтому $\rho^n\ge1/\delta(\Lambda)\ge1/\delta_{\max}(n)$. По
теореме Кабатянского--Левенштейна $\delta_{\max}(n)\le2^{-(0{,}599+o(1))n}$, и
значит
\[
\rho\ \ge\ 2^{0{,}599}-o(1)\ =\ 1{,}5147-o(1).
\]
Эйзенштейнову шагу нужно $\rho\le1{,}52753$: он живёт в окне шириной
\textbf{0{,}85\,\%} над абсолютным препятствием. Это и объясняет, почему
$7^{n/2}$ существует лишь в исключительных размерностях $n=2,4,6,8,24$, где
$\rho=\sqrt2$ или меньше, и не может быть общим асимптотическим механизмом ---
тогда как $\rho\le2$ достижимо в любой размерности, что и есть $3^n$
Лармана--Роджерса. Отрицательные результаты работы для $K_{12}$
($\rho=1{,}633$) и для эйзенштейнова ранга~5 ($\rho=1{,}764$) --- частные
проявления той же тесноты.
